\section{Introduction}
\label{sec:introduction}

Agent Skills turn instructions into supply-chain artifacts. A Skill can bundle
persistent natural-language directives with scripts, manifests, installation
commands, and external resources. Once admitted, this package influences an
agent while inheriting its access to files, credentials, networks, and
processes. The first security decision is therefore not whether a particular
action should run, but whether the package should enter the agent's trust
boundary at all.

The threat is already visible in public ecosystems. A recent measurement of
98,380 Skills confirmed 157 malicious packages through static triage, sandbox
execution, and expert review~\cite{liu2026maliciousskills}. This workflow
produces strong behavioral evidence, but only after a candidate is executed and
its trigger is reached. Runtime defenses likewise mediate an agent after
admission. Registries, enterprise gateways, and local agents still need a
pre-installation filter that can inspect untrusted packages without invoking
their code or dependencies.

That constraint makes the problem fundamentally ambiguous. Sensitive file
access, network transfer, process creation, and credential use are both useful
automation primitives and components of collection, exfiltration, or remote
execution. Instructions may state intent, conceal it, or redefine the apparent
purpose of the package. Static paths connect operations across files, yet a
connected path alone does not establish authorization. Existing Skill
detectors recover increasingly rich artifact evidence~\cite{wang2026malskills},
but recent source-disjoint evaluation shows that repository provenance can
dominate apparent accuracy~\cite{wang2026maliciousskillbench}. The unresolved
question is therefore empirical: \emph{does structured, source-grounded static
evidence improve malicious-Skill detection beyond direct semantic review when
provenance is separated?}

We answer this question with \emph{Read Before You Run}, a zero-execution
analysis pipeline. The detector reads every target file as bounded data and
keeps the package's boundary description separate from operational
instructions and implementation. It normalizes sensitive objects and
operations, constructs cross-file Python call graphs with fixed-point
interprocedural taint summaries, and preserves unresolved calls, unsupported
languages, truncation, and invalid outputs. A schema-constrained reviewer can
therefore return \textsc{Pass}, \textsc{Review}, or \textsc{Block} while citing
the evidence used for each finding.

The source-disjoint experiment rejects the hoped-for accuracy advantage. On
the 977 packages successfully processed by all compared methods, Ours produces
one false positive but detects only 57/491 malicious Skills. A one-call model
that sees only the primary Skill document detects 109/491 at comparable
precision and is significantly more accurate. The independently extracted
high-level function also provides no measurable benefit on the smaller
benchmark. In contrast, the same pipeline appears strong on a 200-package set
whose labels are perfectly confounded with provenance. The comparison exposes
a practical boundary: structured evidence supports low-false-positive triage
and auditable inspection, but does not yet improve source-disjoint detection
over direct document review.

This paper makes three contributions:
\begin{itemize}
  \item \textbf{A traceable zero-execution pipeline} that connects constrained
  instruction analysis, versioned sensitive-object rules, cross-file call
  recovery, and interprocedural taint summaries without loading target code.
  \item \textbf{A source-disjoint comparison} against frozen rules and direct
  document classification, reporting paired significance, processing failures,
  model cost, and three-way review workload rather than accuracy alone.
  \item \textbf{A falsification result for structured Skill detection:} the
  complete pipeline and its high-level function view fail to improve detection
  over direct semantic review, while retaining a low-FPR, evidence-bearing
  operating point and scalable audit front end.
\end{itemize}

Section~\ref{sec:background} defines the threat model,
Section~\ref{sec:related} reviews adjacent defenses,
Section~\ref{sec:methodology} presents the detector, and
Section~\ref{sec:evaluation} evaluates its supported and rejected claims.
